%%=============================================================
%% Capítulo 6 — Conclusiones y trabajo futuro
%%=============================================================
\chapter{Conclusiones}
\label{chap:conclusiones}

Este trabajo ha abordado el diseño y la implementación parcial de un sistema de
inteligencia artificial para la automatización parcial del protocolo \ac{PAALSS},
orientado a reducir la carga de trabajo del especialista en el análisis de muestras
de lenguaje producidas por usuarios de \ac{SAAC} en español. El punto de partida era
un proceso íntegramente manual que consume entre tres y ocho horas de trabajo experto
por cada sesión de treinta minutos~\cite{soto2020protocolo}; el punto de llegada es un
\emph{pipeline} modular que automatiza las etapas más mecánicas de ese proceso y las
articula en una interfaz local que el clínico puede operar sin conocimientos de
programación.

En los apartados que siguen se describe de forma separada lo que ha quedado
\textbf{implementado y probado}, lo que ha quedado \textbf{diseñado pero pendiente de
implementación}, y las \textbf{líneas de trabajo futuro} que permitirían completar y
consolidar el sistema, incluyendo la validación formal frente a anotaciones expertas.


%%=============================================================
\section{Trabajo realizado por objetivo}
\label{sec:conclusiones_oe}
%%=============================================================

La Tabla~\ref{tab:cumplimiento} resume el grado de cumplimiento de cada objetivo
específico del trabajo.

\begin{table}[ht]
  \footnotesize
  \renewcommand{\arraystretch}{1.3}
  \setlength{\tabcolsep}{4pt}
  \begin{center}
  \begin{tabular}{|L{1.2cm}|L{3.5cm}|L{3.5cm}|L{5.2cm}|}
  \hline
  \rowcolor{tema!10}
  \bft{Obj.} & \bft{Descripción} & \bft{Estado} & \bft{Observaciones} \\ \hline
  OE1 & Extracción de \ac{CT} desde audio
      & \emph{Implementado y probado (canal auditivo)}
      & ASR + diarización + detección de eco + PANNs; canal vídeo pendiente \\ \hline
  OE2 & Pipeline de procesamiento lingüístico
      & \emph{Implementado (reglas + spaCy)}
      & Vocabulario, morfología básica, \ac{LME} y \ac{SVO}; pendiente validación \\ \hline
  OE3 & Clasificador morfológico \emph{few-shot}
      & \emph{Implementado (reglas + BETO \emph{zero-shot})}
      & Infinitivos, participios y número; BETO opcional (lazy load) \\ \hline
  OE4 & Pre-rellenado \ac{RAG} del formulario
      & \emph{Parcialmente implementado}
      & FAISS + métricas operativos; integración Ollama pendiente de prueba \\ \hline
  \end{tabular}
  \end{center}
  \caption[Cumplimiento de objetivos específicos]{Grado de cumplimiento de los objetivos
    específicos al cierre de este trabajo.}
  \label{tab:cumplimiento}
  \renewcommand{\arraystretch}{1}
\end{table}

\subsection{OE1 — Extracción automática de turnos comunicativos}
\label{subsec:conc_oe1}

El objetivo de extracción de \ac{CT} a partir de la grabación de vídeo está
implementado en su canal de audio. La cadena de procesamiento es funcional de extremo
a extremo: extracción del audio mediante \texttt{ffmpeg}, transcripción con
\emph{Whisper} \texttt{large-v3}~\cite{whisper_radford2023}, diarización con
\emph{pyannote.audio}~\cite{pyannote_plaquet2023} con número fijo de hablantes, fusión
mediante el enfoque DiCoW~\cite{dicow_diarization}, recuperación de enunciados del
comunicador por eco y clasificación de eventos no verbales mediante
\emph{PANNs}~\cite{pyannote_plaquet2023}. El resultado se almacena en el contrato
de datos estándar \texttt{*\_cts\_paalss.json}, definido con modelos \emph{Pydantic v2}
en \texttt{src/cts\_schema.py}.

El canal de vídeo, necesario para recuperar los pictogramas seleccionados en el
comunicador sin activación de síntesis de voz, está diseñado en la arquitectura
general del sistema (véase la Figura~\ref{fig:arquitectura}) pero queda fuera del
alcance de este prototipo. Un análisis de cobertura sobre la sesión~1 muestra que el
canal de audio, incluso combinando los tres módulos de recuperación, tiene un techo
teórico de aproximadamente el 47\,\% de los turnos comunicativos de la aprendiz,
lo que confirma que el canal de vídeo es una extensión obligatoria para alcanzar una
cobertura clínicamente útil.

\subsection{OE2 y OE3 — Procesamiento lingüístico y clasificación morfológica}
\label{subsec:conc_oe2_oe3}

El módulo de análisis lingüístico \texttt{src/pipeline/etapa3\_nlp.py} está
implementado. Aplica el modelo \texttt{es\_core\_news\_lg} de \emph{spaCy} sobre el
texto de cada \ac{CT} válido para calcular la distribución de categorías
gramaticales, los rasgos morfológicos básicos, la \ac{LME} en morfemas y la presencia
del patrón \ac{SVO}. Los resultados se almacenan como instancias de
\texttt{LinguisticAnalysis} en el campo \texttt{analisis} del esquema \ac{JSON}.

El clasificador morfológico para las categorías ambiguas del lenguaje asistido está
implementado en \texttt{src/pipeline/etapa3\_morfologia.py} mediante reglas lingüísticas
y confirmación opcional con \emph{BETO}~\cite{beto_morphological} \emph{zero-shot}
(\ac{PET}~\cite{fewshot_clinical_nlp}). La evaluación formal de la exactitud de este
módulo frente a las anotaciones expertas queda pendiente como trabajo futuro, junto
con la evaluación completa del sistema.

\subsection{OE4 — Pre-rellenado automático del formulario}
\label{subsec:conc_oe4}

El módulo de pre-rellenado del formulario \ac{PAALSS} ha sido diseñado en su
arquitectura completa (\ac{RAG} + \emph{FAISS} + \emph{Ollama}) y parcialmente
implementado. Los componentes de cálculo de métricas (\texttt{aggregate\_metrics}),
vectorización y construcción del índice \emph{FAISS} están operativos. La integración
con el modelo de lenguaje local mediante \emph{Ollama} para la generación de
propuestas campo a campo del formulario está implementada en el módulo
\texttt{src/pipeline/etapa4\_formulario.py} pero no ha sido probada de extremo a extremo
por ausencia de las métricas \ac{NLP} completas que requiere (resultado de OE2).
El relleno estructurado del \ac{PDF} \emph{AcroForm} de Gloria Soto está implementado
y disponible en la interfaz \emph{Streamlit}.

%%=============================================================
\section{Limitaciones del trabajo}
\label{sec:conclusiones_limitaciones}
%%=============================================================

Las limitaciones más relevantes del trabajo pueden clasificarse en tres categorías.

La primera es estructural y no eliminable: el canal de audio tiene un techo de
cobertura del 47\,\% de los \ac{CT} de la aprendiz, determinado por la proporción de
turnos producidos mediante selección de pictogramas sin síntesis de voz. Ninguna
mejora en los modelos de transcripción o diarización puede superar este techo; el
único camino es incorporar el canal de vídeo.

La segunda es de alcance: el sistema procesa únicamente español y sesiones con
formato de dos participantes grabados en vídeo con audio inteligible, condiciones
que coinciden con el diseño del \ac{PAALSS} pero que excluyen un gran número de
situaciones clínicas reales. La restricción de procesamiento local, si bien es un
requisito de privacidad no negociable, limita el tamaño de los modelos de lenguaje
utilizables y, por tanto, la calidad de las propuestas generadas por el módulo \ac{RAG}.

La tercera es de datos: el \emph{corpus} de referencia consta de dos sesiones de una
única aprendiz, con un único interlocutor y un único tipo de actividad (lectura
compartida). Los resultados del prototipo, cuando se evalúen formalmente, no podrán
generalizarse con confianza a otras usuarias, otros comunicadores ni otros contextos
de interacción sin una ampliación significativa del corpus.

La cuarta es de evaluación: ninguna de las métricas de comparación frente al estándar
experto ha sido calculada. Los valores numéricos presentados en el capítulo anterior
(TNW, NDW, \ac{LME}, distribución de categorías gramaticales) son salidas del análisis
automático que no han sido contrastadas sistemáticamente con las anotaciones de la
especialista. La evaluación formal queda como la prioridad más urgente del trabajo futuro.


%%=============================================================
\section{Trabajo futuro}
\label{sec:conclusiones_futuro}
%%=============================================================

Las líneas de trabajo futuro se ordenan a continuación por su impacto esperado sobre
la utilidad clínica del sistema.

La tarea más inmediata y prioritaria es la \textbf{validación formal del sistema}
frente a las anotaciones expertas disponibles. El protocolo contempla la evaluación
modular de cada componente: transcripción mediante \ac{WER}, diarización mediante
\ac{DER}, morfología y sintaxis mediante precisión, exhaustividad y F1 por categoría,
y pre-rellenado mediante \emph{kappa de Cohen} campo a campo. Esta validación requiere
completar previamente la integración de \emph{Ollama} en el módulo \ac{RAG} (OE4).

En segundo lugar, la \textbf{implementación del canal de vídeo} para la detección de
pictogramas seleccionados en el comunicador es la única vía para superar el techo del
47\,\% del canal de audio. Esta extensión requiere la detección de la región de interés
del tablero en el frame, la clasificación de los pictogramas activados y su fusión
temporal con los \ac{CT} del canal de audio.

En tercer lugar, la \textbf{ampliación del corpus de evaluación} a nuevas sesiones y
nuevas aprendices, en colaboración con la \emph{Fábrica de Palabras}, permitiría
estimar la generalización del sistema y detectar los módulos con mayor impacto en la
variabilidad inter-usuario.

Por último, el \textbf{empaquetamiento del sistema como instalador autónomo} —que
incluya la descarga y configuración automática de los modelos, los entornos de
ejecución y la interfaz, de modo que el clínico pueda instalar y ejecutar el
prototipo sin asistencia técnica— es una condición necesaria para la realización de
la demostración con profesionales de la \emph{Fábrica de Palabras} y para cualquier
uso clínico más allá del entorno de desarrollo.


%%=============================================================
\section{Reflexión final}
\label{sec:conclusiones_reflexion}
%%=============================================================

Este trabajo ha demostrado que es técnicamente viable construir, con modelos de
código abierto y procesamiento completamente local, una cadena de análisis que va
desde el audio crudo de una sesión de comunicación asistida hasta un conjunto de
métricas lingüísticas estructuradas y una propuesta de formulario \ac{PAALSS}
parcialmente cumplimentado. La principal aportación no es un sistema acabado, sino
un \emph{pipeline} modular documentado, con contratos de datos explícitos y una
arquitectura que permite sustituir o mejorar cada componente de forma independiente
a medida que el corpus de evaluación crezca y las limitaciones actuales se superen.

El análisis de cobertura del canal de audio —que sitúa el techo teórico en el
47\,\%— es en sí mismo un resultado de valor para la comunidad investigadora, ya
que cuantifica por primera vez, sobre datos reales de una sesión \ac{PAALSS}, la
proporción de turnos comunicativos de una usuaria de \ac{SAAC} que son recuperables
desde el canal auditivo. Este resultado valida empíricamente la necesidad del canal
de vídeo como componente obligatorio de cualquier sistema de análisis automático de
lenguaje asistido, más allá de las justificaciones teóricas disponibles en la literatura.

Del mismo modo, la elección de modelos locales y la definición de un protocolo de
evaluación frente a anotaciones expertas establecen las bases para que este trabajo
pueda progresar hasta una demostración clínica real con la \emph{Fábrica de Palabras},
contribuyendo así al objetivo más amplio de hacer el \ac{PAALSS} accesible en entornos
clínicos y educativos donde la carga de trabajo manual impide su aplicación sistemática.


% Local Variables:
%  coding: utf-8
%  mode: latex
%  mode: flyspell
%  ispell-local-dictionary: "castellano8"
% End:
